\documentclass[12pt]{article}

\usepackage[T1]{fontenc}
\usepackage[utf8]{inputenc}
\usepackage[brazil]{babel}
\usepackage{lmodern}
\usepackage{microtype}
\usepackage[a4paper,margin=2.5cm]{geometry}
\usepackage{setspace}
\usepackage{indentfirst}

\setstretch{1.15}

\title{Economia Keynesiana, Nova Keynesiana e Nova Clássica \thanks{Texto Traduzido livremente por Luiz Mario Andrade com o auxílio da ferramenta Chat GPT 5.4}}
\author{Bruce Greenwald e Joseph E. Stiglitz}
\date{Fevereiro de 1987}

\begin{document}

\maketitle

\begin{center}
    SÉRIE DE ARTIGOS DE TRABALHO DO NBER \\[1cm]
    Trabalho de Pesquisa No. 2160 \\[1cm]
    NATIONAL BUREAU OF ECONOMIC RESEARCH \\
    1050 Massachusetts Avenue \\
    Cambridge, MA 02138 \\
    Fevereiro de 1987
\end{center}

\vspace{1cm}

A pesquisa aqui relatada faz parte do programa de pesquisa do NBER em Flutuações Econômicas. Quaisquer opiniões expressas são as dos autores e não necessariamente as do National Bureau of Economic Research.

\newpage

\begin{abstract}
Grande parte da nova teoria da macroeconomia que foi construída sobre modelos microeconômicos de informação imperfeita leva a conclusões que são surpreendentemente próximas em espírito à análise original de Keynes. Este artigo resume as implicações macroeconômicas de modelos baseados em informação de salários de eficiência, racionamento de crédito e a quebra dos mercados financeiros para títulos do tipo capital próprio (ações). Ele mostra como esses modelos levam a comportamentos das empresas e interações entre agentes econômicos que explicam muitos dos fenômenos identificados por Keynes em termos qualitativos, os quais foram amplamente perdidos em formalizações subsequentes do modelo keynesiano. Esses macromodelos de informação imperfeita fornecem explicações teóricas consistentes, no espírito keynesiano, para o desemprego, ciclos econômicos concentrados em investimentos, preços rígidos e a eficácia das intervenções de política monetária e fiscal. Ao fazê-lo, eles reconciliam a análise macro e microeconômica de uma forma que até agora não foi alcançada nem pelos keynesianos tradicionais, que ignoraram a dimensão micro do problema, nem pelos novos economistas clássicos, que ignoraram a dimensão macro do problema.
\end{abstract}

\vspace{2cm}

\noindent
\textbf{Bruce Greenwald} \\
Bell Communications Research \\
Economic Research Group \\
Room 2A-347 \\
435 South Street \\
Morristown, NJ 07960—1961

\vspace{0.5cm}

\noindent
\textbf{Joseph E. Stiglitz} \\
Department of Economics \\
Dickinson Hall \\
Princeton University \\
Princeton, NJ 08544

\newpage

\section*{Economia Keynesiana, Nova Keynesiana e Nova Clássica}
\begin{center}
    B. Greenwald e J. E. Stiglitz\footnote{O apoio financeiro da National Science Foundation e da Hoover Institution é gentilmente reconhecido. Os comentários de B. são gentilmente agradecidos.}
\end{center}

Por mais de dois séculos, houve duas visões opostas da economia capitalista. Uma, que geralmente atribui suas origens a Adam Smith, enfatiza a eficiência da economia de mercado, a capacidade do sistema de preços de transmitir informações vitais de produtores para consumidores, e vice-versa, e de coordenar decisões de alocação de uma maneira muito além da capacidade de qualquer planejador central. A outra concentrou sua preocupação nas falhas do capitalismo, particularmente nos episódios periódicos de desemprego massivo de capital e trabalho. Certamente, afirmam os adeptos desta visão, estes não podem ser manifestações de um sistema econômico eficiente.

Para os crentes clássicos na eficiência das economias de mercado, esses episódios eram vistos como situações de desequilíbrio, aberrações temporárias de uma economia de outra forma eficiente; as forças de mercado, se deixadas por si mesmas, restaurariam rapidamente o equilíbrio. Seus descendentes modernos, os Novos Economistas Clássicos, foram um passo além: eles negam a própria existência de um problema; as mudanças massivas nos níveis de emprego podem ser melhor interpretadas, em sua visão, como uma resposta racional a mudanças nos preços relativos (trabalhadores optaram por ter mais lazer em 1932 por causa do baixo salário relativo). O desemprego de capital não seria mais sério do que o fato de o pneu sobressalente de um carro estar, na maior parte do tempo, sem uso: a capacidade extra é mantida para aqueles poucos momentos em que é realmente necessária.

Para Keynes e seus seguidores modernos, essas visões não são apenas absurdas: elas ridicularizam o "método científico" que seus adeptos afirmam seguir. Pior ainda, são irresponsáveis: na medida em que os governos seguem as políticas de não intervenção frequentemente preconizadas, eles não apenas condenam aqueles indivíduos que não conseguem obter emprego remunerado à privação econômica resultante, mas também condenam a sociedade que tolera esse desemprego a uma série de consequências sociais e econômicas que decorrem desse desemprego.

Uma das grandes contribuições de Keynes foi, na verdade, uma reconciliação das duas visões opostas do capitalismo: em vez de negar a existência do problema do desemprego ou sua importância, ele o enfrentou de frente, argumentou que a intervenção governamental limitada poderia corrigir esse mal e, com esse mal corrigido, a economia voltaria a operar de maneira eficiente: a visão clássica seria então restaurada. Samuelson apelidou isso de síntese neoclássica.

As próprias razões do sucesso da abordagem de Keynes forneceram a base para a eventual desilusão que se instalou tão fortemente, pelo menos nos Estados Unidos, ao longo das últimas duas décadas. Keynes tentou reter o máximo possível do aparato clássico (neoclássico); o modelo padrão foi alterado de formas mínimas, com consequências dramáticas. A síntese neoclássica, tão atraente ideologicamente quanto era para aqueles que acreditavam no sistema de mercado, mas estavam perturbados pelo desemprego massivo, foi aceita como um artigo de fé, não derivada de nenhuma estrutura teórica geral: a questão fundamental de por que as falhas da economia de mercado deveriam ocorrer apenas nas doses massivas que associamos às suas recessões e depressões periódicas nunca foi feita, muito menos respondida. Não era mais plausível que a Grande Depressão fosse apenas a pior manifestação de um conjunto de ineficiências que, embora generalizadas, eram muito mais difíceis de detectar ou provar?

A esquizofrenia a que a economia keynesiana deu origem refletiu-se na maneira como a economia era ensinada: cursos de microeconomia, nos quais os alunos eram apresentados à mão invisível de Adam Smith e aos teoremas fundamentais da economia do bem-estar, eram seguidos por cursos de macroeconomia, focando nas falhas da economia de mercado e no papel do governo em corrigi-las. Duas subdisciplinas se desenvolveram, com os microeconomistas menosprezando a (falta de) rigor dos macroeconomistas e denegrindo a falta de fundamentos teóricos, enquanto os macroeconomistas castigavam os microeconomistas pela óbvia inadequação de suas teorias.

A insatisfação com a economia keynesiana baseava-se não apenas nessa esquizofrenia, mas também na falta de justificativa para algumas das suposições centrais, por exemplo, as rigidezes de salários e preços. Salários e preços não eram rígidos. Se eles não caíssem o suficiente, por que não caíam? Por que as empresas que queriam vender mais simplesmente não baixavam seus preços? Um quarto de século de pesquisa não conseguiu fornecer respostas convincentes para essas perguntas.\footnote{Estas não foram as únicas objeções à teoria keynesiana, nem as únicas fontes de insatisfação. Historicamente, a incapacidade de lidar com o surto simultâneo de inflação e desemprego no início dos anos setenta pode ter desempenhado um papel tão crítico quanto os problemas teóricos mais profundos com os quais nos preocupamos aqui. O crescimento do monetarismo durante este período pode ser devido mais à sua capacidade de fornecer prescrições simples e claras do que à sua capacidade de remediar as deficiências teóricas na análise keynesiana. Os teóricos estavam preocupados, é claro, não apenas com sua incompletude e inconsistência com a análise microeconômica tradicional, mas também com suas inconsistências internas, por exemplo, suas suposições de expectativas inconsistentes.} Este estado de coisas não poderia continuar por muito tempo. Uma tentativa de reconciliação entre os dois parecia inevitável.

Havia, de modo geral, duas maneiras pelas quais as abordagens alternativas poderiam ser reconciliadas: adaptar a macroteoria à microteoria, ou o inverso. A Nova Economia Clássica seguiu a primeira abordagem. Argumentou que o que estava errado com a macroeconomia era a sua ausência de fundamentos microeconômicos rigorosos. Seus defensores iniciaram um ambicioso programa de pesquisa, visando derivar o comportamento dinâmico e agregado da economia a partir dos princípios básicos de empresas e indivíduos racionais e maximizadores. A Escola reconheceu a importância da dinâmica para a compreensão do macrocomportamento e reconheceu o papel central das expectativas na determinação do comportamento dinâmico. Concentrou sua atenção, então, nas consequências da formação de expectativas racionais, e é este aspecto de seu trabalho que deu à Escola seu nome alternativo.

O nome escola de "Expectativas Racionais" é, no entanto, enganoso: as doutrinas centrais da abordagem derivam não de sua crença em expectativas racionais, por mais plausível ou implausível que essa suposição possa ser; mas sim de suas antigas suposições clássicas de equilíbrio de mercado. E com essas suposições, a conclusão de que não há desemprego e a irrelevância da macropolítica governamental\footnote{De fato, Neary e Stiglitz (1983) mostraram que, com expectativas racionais e rigidezes de preços, a política governamental é ainda mais eficaz do que com expectativas míopes: os multiplicadores são ainda maiores. Nos modelos Novos Clássicos, a macropolítica é irrelevante porque é desnecessária. Mesmo assumindo o pleno emprego, as políticas governamentais terão, em geral, efeitos reais, por exemplo, na acumulação de capital. Veja Stiglitz (1981, 1983).} seguem como consequências triviais. Essas conclusões, deve-se notar, seguiriam mesmo se os participantes do mercado fossem muito menos racionais do que o postulado pela teoria.

A outra abordagem vê o desemprego, o racionamento de crédito, os ciclos econômicos como problemas econômicos reais, fenômenos que não podem ser reconciliados com a microteoria padrão e, portanto, busca como seu objetivo o desenvolvimento de uma microteoria que possa explicar esses fenômenos. Por falta de um termo melhor, referir-me-ei a isso como a Nova Economia Keynesiana. O trabalho nesta área centrou-se na compreensão das consequências da informação imperfeita e dos mercados incompletos, tanto para a microeconomia quanto para a macroeconomia. Como a Nova Economia Clássica, ela busca uma teoria única, mas, ao contrário da Nova Economia Clássica, busca explicar o desemprego, em vez de negar sua existência. E o desemprego mostra-se apenas uma manifestação de um conjunto muito mais amplo de falhas de mercado.

Este artigo tenta apresentar os contornos gerais da Nova Economia Keynesiana e mostrar em que aspectos ela é semelhante à economia keynesiana tradicional e em que aspectos difere. Keynes tinha uma visão de como a economia funcionava que era marcadamente diferente da teoria neoclássica padrão. As decisões das empresas não eram baseadas em cálculos racionais. Keynes usou uma linguagem pitoresca para descrever o comportamento dos empreendedores: eles eram movidos por "espíritos animais". Mas quando Keynes se propôs a escrever um modelo simples, ele recorreu a modos de pensamento mais tradicionais, e esses aspectos foram reforçados nos desenvolvimentos subsequentes (por exemplo, por Hicks). É nossa tese que os problemas de Keynes surgiram de sua incapacidade de romper com seu treinamento neoclássico, que sua visão, capturada tão bem em muitas de suas passagens brilhantemente escritas, fornece maior percepção para entender o desemprego e os ciclos econômicos do que o modelo keynesiano formal.

\section*{Algumas Percepções Keynesianas Fundamentais}

Existem quatro percepções de Keynes que consideramos essenciais na construção de um modelo de desemprego e flutuações econômicas:

1. Uma teoria geral deve explicar a persistência do desemprego, bem como a ciclicidade em certas variáveis econômicas fundamentais. Para explicar a persistência do desemprego, é necessário desenvolver uma teoria do mercado de trabalho. Keynes recorreu, neste ponto, à suposição de rigidez salarial (real ou nominal?), uma suposição que foi colocada no centro da escola de preços fixos, mas que foi atacada por críticos da teoria keynesiana tanto em bases empíricas — os salários, afinal, caíram um terço na Grande Depressão, e países enfrentando inflação, presumivelmente com algum espaço para mudanças no salário real, enfrentaram desemprego tal como países nos quais salários e preços caíram — quanto em bases teóricas — nenhuma explicação para rigidezes salariais, além de considerações institucionais ad hoc, é fornecida.

O que é necessário, contudo, para as conclusões keynesianas, não são salários absolutamente rígidos, mas apenas que os salários não caiam para os níveis de equilíbrio de mercado.\footnote{Uma abordagem alternativa à discutida aqui é que, embora os salários caiam, os preços caem aproximadamente na mesma proporção. Assim, os salários reais não caem. (Esta é a abordagem adotada por Solow-Stiglitz (1968)). Embora a queda de salários e preços possa dar origem a um aumento na demanda por bens de consumo — o efeito Pigou-Patinkin (na ausência de considerações de Barro-Ricardo) —, há pouca dúvida de que, no curto prazo, no período para o qual a análise keynesiana é apropriada, este efeito não tem significância quantitativa. Deve-se notar que em muitos modelos macroeconômicos, os efeitos de saldo real desempenham um papel crucial, se um tanto sub-reptício. Assim, no modelo de preço fixo, o desemprego é atribuído a um nível excessivamente alto de salários e preços; se salários e preços caíssem (mantendo os salários reais, digamos, constantes), o emprego aumentaria; mas isso ocorre apenas por causa de um suposto efeito de saldo real importante. Deve-se notar também que, no curto prazo, mesmo que houvesse efeitos de saldo real do tipo observado por Pigou e Patinkin, esses efeitos poderiam ser compensados, no curto prazo, por efeitos dinâmicos de substituição intertemporal; se os preços caem e os consumidores acreditam que cairão ainda mais, isso pode reduzir sua demanda atual por consumo. Veja Neary e Stiglitz (1983) ou Grandmont.} As teorias de salários de eficiência\footnote{Para uma pesquisa, veja Stiglitz (1986a, 1986b).} recentemente desenvolvidas fornecem tal explicação. Baseiam-se na hipótese de que existe informação imperfeita sobre as características dos trabalhadores; que as ações dos trabalhadores não podem ser monitoradas; e que não é possível redigir contratos que garantam que o trabalhador arque com todas as consequências de suas ações. Como resultado, a qualidade da força de trabalho, sua produtividade (e, portanto, os lucros da empresa) pode aumentar com o salário pago. Da mesma forma, a rotatividade de pessoal (turnover) pode diminuir com um aumento no salário e, como a empresa deve arcar com parte dos custos de rotatividade, novamente os lucros podem aumentar com um aumento nos salários, até certo ponto. Diante do desemprego, os salários podem não cair, pois as empresas reconhecerão que, se reduzirem os salários, a produtividade diminuirá, a rotatividade poderá aumentar e os lucros cairão. Nessa perspectiva, as empresas são competitivas; existem muitas empresas no mercado; mas, no entanto, as empresas são definidoras de salários (wage setters), pelo menos dentro de um intervalo. Se o salário walrasiano, onde a demanda por trabalho é igual à oferta, for muito baixo, qualquer empresa tem a opção de aumentar seu salário e, assim, aumentar seus lucros. O salário de eficiência, o salário que maximiza os lucros da empresa, pode, é claro, variar com as circunstâncias econômicas; portanto, o salário não é absolutamente rígido. Mas os salários não precisam cair para os níveis de limpeza de mercado.\footnote{Assim, as implicações de política dessas teorias podem ser marcadamente diferentes daquelas dos modelos padrão de salários e preços fixos. Estes últimos assumem que a política econômica não tem efeito sobre os salários pagos. Os modelos de salários de eficiência reconhecem que certas políticas (por exemplo, compensação de desemprego) podem ter efeitos fortes nos salários de equilíbrio, e as consequências disso precisam ser levadas em conta.}

Uma objeção a esta teoria — bem como à teoria keynesiana padrão — é que a presença de rigidezes salariais em alguns setor(es) da economia não é suficiente para explicar o desemprego.\footnote{Esta não é, obviamente, a única objeção à teoria dos salários de eficiência. Para uma discussão mais extensa, veja Stiglitz (1986b).} Enquanto houver algum setor com salários flexíveis, qualquer indivíduo que opte por não trabalhar lá está voluntariamente desempregado. Consideramos que esta é uma objeção amplamente semântica: o fato é que indivíduos que são observacionalmente indistinguíveis do indivíduo desempregado estão sendo empregados com salários mais altos; que o equilíbrio de mercado é ineficiente; e que recursos que poderiam ser empregados produtivamente permanecem ociosos. (Em outro lugar, Greenwald-Stiglitz (1986b), discutimos extensamente uma variedade de razões pelas quais pode ser racional para um indivíduo não aceitar um salário baixo atualmente, se ele acredita que um emprego melhor remunerado se tornará disponível em um futuro próximo; estas têm a ver com informação assimétrica, com a informação transmitida pela disposição do indivíduo em aceitar um emprego de baixo salário, bem como com o fato de que, uma vez empregado, ele se torna "trabalho usado" com efeitos adversos sobre salários futuros, semelhantes aos que surgem no modelo de limões de Akerlof. Também discutimos as razões pelas quais um trabalhador pode não estar disposto a aceitar um salário baixo de um empregador, com a promessa de um salário futuro mais alto se a empresa sobreviver, porque fazê-lo tornaria, de fato, o trabalhador um detentor de capital próprio na empresa; veja abaixo uma discussão sobre por que os indivíduos não desejariam fazê-lo.)

2. Uma teoria macroeconômica deve não apenas explicar a persistência do desemprego, mas também suas flutuações. Esta não é uma tarefa tão fácil quanto pode parecer. Existem dois problemas.

Primeiro, quais são as fontes de choques na economia que dão origem a tão grandes perturbações na atividade econômica?\footnote{O problema não é que não existam choques, mas a magnitude dos choques e suas correlações entre setores.} É difícil encontrar choques que sejam externos ao sistema econômico, embora ocasionalmente tais choques — como guerras e o aumento nos preços do petróleo\footnote{Mas mesmo isso deve ser visto como endógeno: nenhuma mudança na demanda ou na oferta ocorreu, embora a formação de um cartel possa ser vista, sob essa perspectiva, como um choque exógeno ao sistema.} — ocorram. E se não houver choques exógenos, quais são os mecanismos internos que levam a tais variações na atividade econômica?\footnote{A crença de que o sistema econômico capitalista não poderia ser internamente falho ao ponto de dar origem endogenamente a flutuações fornece grande parte da motivação para a busca de um culpado externo, geralmente identificado como o "governo" ou as "autoridades monetárias".} Há evidências consideráveis de que grande parte das flutuações se deve a variações na demanda por investimento e, em particular, por estoques. Mas, em teoria, com funções de produção côncavas, sendo o preço sombra do trabalho baixo nas recessões e com taxas de juros reais baixas, deveria haver suavização da produção: os estoques deveriam servir para amortecer as flutuações econômicas; eles não deveriam servir para exacerbá-las.\footnote{Existem outros estabilizadores na economia que não discutimos. A poupança serve para estabilizar o consumo. Na ausência de poupança adequada, o seguro fornecido por contratos implícitos serve para estabilizar rendas e, assim, para estabilizar o consumo. Assim, os contratos implícitos, em vez de exacerbar as flutuações econômicas, podem realmente servir para reduzi-las. (Para uma discussão mais extensa deste artigo, veja Stiglitz (1986).)}

Keynes estava correto ao enfatizar a importância do investimento para a compreensão das flutuações econômicas. Para explicar as flutuações no investimento, ele teve que confiar nos "espíritos animais", em mudanças inexplicáveis nas expectativas. As suposições sobre expectativas subjacentes à análise de Keynes podem não ser piores do que as subjacentes à escola de expectativas racionais. No entanto, há algo inquietante em confiar tão completamente no inexplicável ou no irracional para uma teoria dos ciclos econômicos.\footnote{Os macroeconomistas de estilo antigo explicavam essas flutuações de investimento por referência a aceleradores. Esses modelos de acelerador baseavam-se implicitamente em empresas extrapolando os níveis atuais de produção para o futuro; isto é, baseavam-se em expectativas irracionais. Esses modelos podem ser descritivamente precisos, mais descritivamente precisos do que os modelos de expectativas racionais; a questão que nos colocamos é, no entanto: pode-se explicar as flutuações no investimento sem recorrer a tais expectativas irracionais? LaRoque apresentou recentemente uma explicação bastante diferente para o comportamento dos estoques. Ele postula expectativas racionais, mas sem limpeza de mercado. Ele argumenta que os estoques são mantidos para fins especulativos e assume que os especuladores têm alguma prioridade sobre outros em caso de excesso de oferta de bens. Ele mostra então que isso pode dar origem a ciclos de estoque. Suspeitamos que, se olhássemos cuidadosamente para os salários sombra durante os períodos recessivos e perguntássemos se, dados esses salários sombra, teria valido a pena as empresas investirem em estoques, armazenarem as mercadorias até um período de alta demanda, levando em conta as taxas de juros de mercado, a resposta seria sim, e de fato os lucros de fazê-lo teriam excedido os lucros especulativos acumulados durante períodos em que os estoques foram realmente acumulados.}

O segundo problema em explicar as flutuações é o seguinte: no modelo microeconômico típico, mudanças nos preços (taxas de juros, salários) servem para amortecer qualquer perturbação na demanda ou na oferta, exatamente como argumentamos anteriormente que os estoques deveriam estabilizar a economia. Assim, mesmo grandes mudanças exógenas (mudanças nas curvas de demanda ou de oferta) podem resultar em pequenas mudanças nos valores de equilíbrio.

Keynes teve que explicar não apenas por que a curva de demanda por investimento se deslocava, mas também por que as mudanças na taxa de juros não podiam compensar os efeitos diretos. Embora sua explicação sobre a fonte do deslocamento na demanda por investimento possa ter sido incompleta, sua explicação sobre por que as mudanças na taxa de juros não podiam compensar esses efeitos não era persuasiva.\footnote{A teoria da armadilha da liquidez foi projetada para explicar um piso na taxa de juros nominal. Mas, enquanto o dinheiro não puder ser portador de juros negativos, resultados semelhantes seguiriam simplesmente da não-negatividade do juro nominal. O que é relevante para o investimento deveria, ostensivamente, ser a taxa de juros real. E o piso nas taxas de juros nominais, combinado com as pressões deflacionárias de uma recessão, dá origem a aumentos na taxa de juros real. Mas as rigidezes de preços, que parecem estar no cerne das interpretações modernas da análise keynesiana, implicariam uma queda nas taxas de juros reais em uma recessão. Como já notamos, os preços caíram, as taxas de juros reais subiram, mas apenas em uma extensão limitada. Em recessões mais recentes, no entanto, os preços têm subido, às vezes a uma taxa suficiente para tornar as taxas de juros reais negativas (e após impostos, taxas de juros reais muito negativas). Essas quedas na taxa de juros real não foram suficientes para restaurar o investimento. (Isso não é surpreendente; as empresas tradicionalmente exigem retornos esperados de vinte a trinta por cento para realizar um projeto de investimento. Uma mudança na taxa de juros real de 2\% para 1\% é pouco mais que um erro de arredondamento.) No final, é a inelasticidade do investimento em relação aos juros que parece ser crucial; isto é, nenhuma mudança plausível nas taxas de juros reais é suficiente para restaurar o investimento. Abaixo, fornecemos uma explicação para isso.} Em sua teoria, a redução na demanda por investimento deveria ter levado a uma queda na taxa de juros real; no entanto, durante a recessão, as taxas de juros reais subiram. E os empresários argumentaram que as taxas de juros tinham pouco efeito em suas decisões de investimento.

Da mesma forma, ele falhou em fornecer uma explicação de por que os salários não mudaram: ele simplesmente afirmou que não mudaram. A teoria do salário de eficiência explica por que os salários podem não cair para os níveis de limpeza de mercado. Uma teoria semelhante do mercado de capitais mostra por que as taxas de juros podem não cair (Veja Stiglitz-Weiss, 1981, 1983, 1985). Mais geralmente, Akerlof e Yellen apontaram que, mesmo quando as empresas deveriam mudar os salários que pagam, elas podem não fazê-lo; eles mostram que a perda de lucros deste comportamento quase racional pode ser pequena, embora a perda para a sociedade possa ser grande. De fato, se as empresas forem avessas ao risco (como argumentaremos abaixo que serão), e se houver alguma incerteza sobre as consequências das mudanças salariais, manter os salários inalterados diante de certas perturbações é totalmente racional. (Novamente, argumentos semelhantes valem para o mercado de capitais.)

Além disso, os modelos de salários de eficiência mostram ainda por que os salários das empresas são interdependentes: o salário ótimo para a empresa $i$ depende dos salários pagos por todas as outras empresas. Esta interdependência pode levar a equilíbrios múltiplos, nos quais nenhuma empresa altera seu salário, mesmo diante de mudanças em sua demanda.\footnote{Novamente, argumentos semelhantes valem para o mercado de capitais e o mercado de produtos. Veja, por exemplo, Stiglitz (1986c).}

Assim, ao explicar as rigidezes de salários, taxas de juros e preços, essas teorias ajudam a explicar por que certas perturbações são multiplicadas pelo sistema econômico, em vez de serem amortecidas.\footnote{Os macroeconomistas, pelo menos desde o artigo básico de Kahn, enfatizaram o papel dos multiplicadores, sem reconhecer quão inconsistentes eles são com a microanálise padrão. A diferença crítica entre economias que amplificam perturbações e aquelas que as amortecem parece ser as rigidezes de preços; nossa análise explica essas rigidezes e, ao fazê-lo, coloca o multiplicador sobre bases teóricas sólidas.}

Na teoria que apresentamos abaixo, existe um conjunto adicional de razões para a "multiplicação" de perturbações. Na presença de mercados incompletos e informação imperfeita, as ações de uma empresa ou indivíduo exercem efeitos do tipo externalidade sobre os outros; a redução da produção por uma empresa, em resposta ao aumento da incerteza ou a uma redução em seu capital de giro, aumenta a incerteza e reduz o capital de giro de outras empresas. Enquanto os ajustes de preços tendem a amortecer as perturbações, os efeitos de externalidade podem (e, nestes casos, o fazem) exacerbá-las.

3. Keynes estava correto ao enfatizar a importância da dicotomia entre poupança e investimento. Com efeito, ele reconheceu que os fundos dentro da empresa eram diferentes dos fundos no setor doméstico. Embora ele sentisse a importância dessa distinção, quando se propôs a escrever seu modelo, ele escreveu o modelo neoclássico padrão, no qual essa distinção não desempenha papel algum. A demanda por investimento era determinada pela taxa de juros. Não havia, por exemplo, racionamento de crédito.

Isso foi ainda mais surpreendente, dada a sua percepção geral (implícita, às vezes explícita) da importância das imperfeições do mercado de capitais. O consumo dos indivíduos estava relacionado à sua renda atual, em parte talvez porque a renda atual era uma boa previsão da renda futura, mas em parte por causa da falta de acesso a fundos. Essa falha em reconhecer suficientemente a importância das imperfeições do mercado de capitais, e as consequências que dela decorrem, acaba sendo o seu erro mais importante.

4. Keynes precisava, como dissemos, encontrar uma fonte de flutuações na atividade econômica. Ficou claro que mudanças na tecnologia, na oferta, não poderiam explicar o que estava ocorrendo na Grande Depressão. Ele, portanto, naturalmente se voltou para mudanças na demanda. Bons economistas formados na tradição marshalliana foram ensinados a separar as perturbações em um mercado em perturbações de demanda e oferta.

Keynes estava correto ao focar nas perturbações de demanda, mas sua dependência do arcabouço marshalliano de demanda/oferta colocou problemas que ele, e seus seguidores, nunca resolveram satisfatoriamente. Pois a teoria marshalliana sugeria que o equilíbrio deveria estar na interseção da demanda e da oferta; se as empresas estivessem em sua curva de oferta, os salários reais dos produtos deveriam subir à medida que o emprego caísse. Esta foi uma das primeiras proposições empíricas da economia keynesiana a cair pelo caminho. Mas, assim como a economia marxista nunca foi abandonada por seus proponentes simplesmente porque suas previsões se revelaram falsas, a economia keynesiana não seria abandonada simplesmente porque uma de suas previsões empíricas importantes se revelou errada. Como sempre na economia, existem três maneiras de lidar com fatos desconfortáveis: (a) negá-los, por exemplo, afirmando que salários e preços são medidos incorretamente (exatamente como os novos economistas clássicos abordam o problema do desemprego, negando a relevância das estatísticas de desemprego); (b) fornecer uma nova interpretação, por exemplo, afirmando que o que é relevante não é o salário à vista (spot wage), devido à existência de contratos de longo prazo (implícitos), ignorando o fato de que os salários reais dos produtos dos trabalhadores recém-contratados ou dos trabalhadores em contratos à vista também não subiram significativamente; (c) afirmar que a proposição empírica não era central para a teoria. Assim, desenvolveu-se uma vasta literatura, afirmando que as empresas, embora resolvessem problemas de maximização intertemporal bastante complicados, agiam como se os preços e as quantidades que enfrentavam fossem fixos.\footnote{Esta literatura remonta aos anos sessenta, com contribuições de Hansen (1951), Solow-Stiglitz (1967) e Barro-Grossman (1971). A literatura subsequente de preços fixos é enorme.} Simplesmente se afirmava que as empresas não usavam a política de preços para afetar as vendas, uma suposição implausível e contrafactual.\footnote{Modelos que postulam empresas imperfeitamente competitivas explicam por que os salários reais podem não ser iguais ao valor do produto marginal; mas eles têm pouco a dizer sobre o desemprego involuntário ou suas flutuações. (De fato, em contraste com os modelos com desemprego clássico, com salários reais em excesso ao valor do produto marginal, aqui os salários reais são menores que o valor do produto marginal; se o emprego é maior ou menor no equilíbrio depende simplesmente das elasticidades da oferta de trabalho (não compensada).) Abaixo, fornecemos uma explicação para a variabilidade cíclica nas margens de lucro (mark-ups). Veja também Stiglitz (1984).}

\section*{A Nova Economia Keynesiana}

A Nova Economia Keynesiana começa com as percepções básicas de Keynes, mas reconhece a dependência excessiva de Keynes em relação a um arcabouço neoclássico e sua falha em reconhecer plenamente as consequências das imperfeições do mercado de capitais, imperfeições que podem ser explicadas em termos de custos de informação.

Os principais ingredientes desta nova perspectiva são os seguintes:

1. O modelo do salário de eficiência, que explica por que os salários, embora não rígidos, não caem para os níveis de limpeza de mercado (veja acima).

2. Imperfeições do mercado de capitais, derivadas da informação imperfeita. Existem assimetrias de informação entre gestores de empresas e investidores potenciais, assimetrias que podem dar origem ao que chamaremos de "racionamento de capital próprio" (equity rationing). O racionamento de capital próprio é importante porque implica que, se as empresas desejam obter mais capital, para investir ou para aumentar a produção, elas devem tomar emprestado os fundos; e mesmo que sejam capazes de fazê-lo, devem se expor a riscos consideráveis, incluindo o risco de falência, ou seja, o risco de não serem capazes de pagar as quantias prometidas.

As consequências disso são exacerbadas pela ausência de mercados de futuros.\footnote{Keynes não comentou explicitamente sobre isso, talvez porque para ele a ausência fosse tão óbvia que não requeria comentário. Para aqueles criados no arcabouço Arrow-Debreu, tornou-se um ritual tentar identificar de que formas o modelo proposto difere daquele paradigma.} Assim, as empresas não podem vender os bens que planejam produzir até que os tenham produzido.\footnote{Existem, é claro, indústrias nas quais os bens são produzidos apenas por encomenda. Isso é particularmente verdadeiro para algumas indústrias de manufatura de bens de investimento. Então, variações no nível de produção devem ser rastreadas até variações nas demandas que elas enfrentam; e estas estão relacionadas à disposição de outros produtores em investir (produzir). Assim, as respostas de oferta de algumas empresas refletem-se nas demandas enfrentadas por outras empresas. Esta é outra razão pela qual a simples dicotomia keynesiana entre perturbações de demanda ou oferta pode ser um tanto enganosa.} Cada decisão de produção é uma decisão de risco, um risco que eles (os gestores e detentores de capital próprio) devem assumir, e que não podem facilmente transferir para outros. A ausência de mercados de futuros implica que as empresas não podem vender sua produção no momento da produção.

Assim, uma análise do comportamento da empresa deve focar em sua disposição de assumir esses riscos. Mudanças inesperadas em sua base de capital de giro (causadas, por exemplo, por mudanças inesperadas nos preços pelos quais ela pode vender seus bens) poderiam, por exemplo, ter um efeito deletério em sua disposição de produzir.

3. Enquanto às vezes isso limita a quantidade que as empresas estão dispostas a produzir, em outros momentos, o acesso das empresas ao capital é limitado; há racionamento de crédito. As razões pelas quais os fornecedores de capital não aumentam as taxas de juros na presença de um excesso de demanda por capital são análogas às razões pelas quais as empresas não baixam os salários na presença de um excesso de oferta de trabalho: o aumento das taxas de juros pode diminuir o retorno esperado para o fornecedor de capital, seja por causa de efeitos de seleção (o mix de candidatos muda adversamente) ou por causa de efeitos de incentivo (os tomadores de empréstimo são induzidos a empreender ações mais arriscadas).

4. A política monetária exerce sua influência — quando o faz — não tanto através da disposição dos indivíduos em manter saldos monetários, mas através da disponibilidade de crédito. Assimetrias de informação implicam que, se os bancos decidirem emprestar menos, não existem outros credores potenciais que sejam substitutos perfeitos. As decisões dos bancos de emprestar são análogas àquelas que determinam a disposição das empresas em produzir. As autoridades monetárias podem tomar ações que afetam a disposição dos bancos em emprestar (ou os termos sob os quais eles estão dispostos a emprestar). Embora, dependendo das circunstâncias econômicas, outros credores possam tomar ações parcialmente compensatórias, suas ações nunca podem ser totalmente compensatórias.

A Nova Economia Keynesiana fornece uma teoria geral da economia, derivada de princípios microeconômicos (e, assim, integra as duas subdisciplinas). Ela consegue tanto preencher as lacunas na teoria keynesiana tradicional (por exemplo, explicando rigidezes salariais parciais, em vez de simplesmente assumir salários rígidos) quanto resolver os paradoxos e inconsistências da teoria keynesiana mais tradicional (tanto as inconsistências internas, por exemplo, sobre como as expectativas são formadas, quanto as inconsistências entre suas previsões e observações). Ela fornece uma explicação tanto para um nível de equilíbrio de desemprego (através das teorias de salários de eficiência) quanto para flutuações econômicas.\footnote{Isso não quer dizer que não existam lacunas importantes na teoria que permanecem. A teoria desenvolvida até agora não fornece um ciclo econômico inteiramente endógeno; ela apenas explica como a economia responde a certos choques. Permanece uma controvérsia sobre se uma teoria de ciclo econômico inteiramente endógena é necessária, ou se deve-se contentar com uma teoria que traduz certos tipos de choques em perturbações nas quais a economia persiste abaixo do "pleno emprego" por um número de períodos. Não tomamos posição sobre essa questão aqui.} A teoria das flutuações econômicas que ela fornece é simples: em linhas gerais, certos choques na economia afetam o estoque de capital de giro das empresas. Mesmo que as empresas tivessem acesso perfeito aos mercados de crédito (ou seja, pudessem tomar emprestado o quanto desejassem, à taxa de juros atuarialmente justa), a quantia que elas estariam dispostas a tomar emprestado é limitada pela sua disposição em assumir riscos; os compromissos fixos associados aos contratos de empréstimo implicam que, à medida que o capital de giro disponível é reduzido, o risco (probabilidade de falência) associado a qualquer nível de empréstimo aumenta. Assim, se o seu capital de giro for reduzido, o seu nível de produção desejado (dado que elas não têm compromissos fixos para vender os seus produtos\footnote{Mesmo que tenham compromissos, os compradores potenciais podem não honrar esses compromissos, particularmente em caso de falência. Nas recessões, o risco associado a qualquer "compromisso" aumenta.}) é reduzido; e leva uma série de períodos antes que os níveis de capital de giro sejam restaurados ao normal. A teoria fornece uma explicação não apenas para por que os choques agregados (como uma diminuição inesperada no nível de preços, resultante de um choque monetário) têm um efeito agregativo na economia, mas também por que os choques setoriais (como uma mudança inesperada na demanda, ou a formação inesperada de um cartel de petróleo) teriam efeitos agregativos: a disposição de produzir será, em geral, uma função côncava do capital de giro e, portanto, uma redistribuição do capital de giro terá efeitos agregativos.\footnote{Esses efeitos de redistribuição parecem ser mais importantes do que os efeitos de redistribuição, por exemplo, às vezes postulados com a política de dívida governamental (a mudança na estrutura de vencimento da dívida tendo um efeito de redistribuição intertemporal ou intrateimporal) ou com algumas formas de seguro. A redistribuição resultante do seguro associado a contratos de trabalho implícitos, uma redistribuição do setor corporativo para o setor doméstico, opera essencialmente através do mecanismo descrito acima. Na presença de mercados de capitais perfeitos, os únicos efeitos decorrentes dessa redistribuição seriam aqueles associados a diferentes propensões marginais a consumir entre capitalistas e trabalhadores.}

Na discussão abaixo, mostraremos como esta teoria fornece uma explicação para vários dos fenômenos que pareciam tão difíceis de explicar para a teoria keynesiana mais tradicional: (a) ela explica por que as empresas não baixam os preços nas recessões, ou seja, explica os movimentos cíclicos nas margens de lucro (mark-ups); (b) ela fornece uma explicação do comportamento cíclico do investimento e dos estoques; (c) ela fornece uma explicação de por que os trabalhadores desempregados não conseguem ser contratados ao se oferecerem para trabalhar por salários mais baixos, e mesmo em indústrias onde considerações de salários de eficiência não são importantes, ela fornece uma explicação parcial de por que os trabalhadores não oferecem trabalhar por salários mais baixos, em troca da promessa de salários mais altos no futuro; e (d) ela fornece uma explicação de por que uma redução imprevista de salários e preços pode, na verdade, servir para exacerbar a recessão, em vez de aliviá-la (ao deteriorar ainda mais a base de capital de giro das empresas).

\section*{Os Erros de Keynes}

Argumentamos que o erro básico de Keynes foi uma dependência excessiva, em sua modelagem formal, das ferramentas neoclássicas/marshallianas que eram, então como agora, o estilo da época. As histórias verbais que acompanham o modelo capturaram muito mais de suas percepções sobre o que estava acontecendo.\footnote{Ele não deve ser julgado com muita severidade: como dissemos, presumivelmente ele desejava tornar suas ideias o mais palatáveis possível para seus contemporâneos e, para fazer isso, ele teve que mostrar que, alterando apenas algumas das premissas básicas do modelo padrão, se poderia obter resultados dramaticamente diferentes.} Pode ser útil nesta conjuntura revisar o que são, em nossa perspectiva, seus erros fundamentais, e sugerir como nossa teoria corrige esses erros. Seus erros mais importantes residem, em nosso julgamento, em sua teoria da empresa e em sua explicação do papel da moeda na determinação do nível de atividade econômica. Ambos podem ser relacionados à sua falha em compreender plenamente a natureza dos mercados de capitais.

1. Keynes falhou em reconhecer a importância da distinção entre títulos de longo prazo e ações (equities). Ele agrupou os dois como ativos de longo prazo. Mesmo na ausência de falência, os dois diferem em suas propriedades de risco, com os títulos subindo de valor nas recessões e as ações caindo. Os dois títulos são, portanto, complementos, em vez de substitutos, nas carteiras dos indivíduos. Para nossos propósitos, no entanto, essa distinção não é tão importante quanto as diferenças na natureza do compromisso das empresas: com títulos e empréstimos, a empresa está comprometida a pagar uma certa quantia em uma data específica; com ações, tal compromisso não existe. Como resultado, tanto para as empresas quanto para os investidores, esses dois títulos estão longe de ser substitutos perfeitos. Particularmente em períodos recessivos, as empresas raramente recorrem ao mercado de ações para levantar o capital necessário: os investidores suspeitam que qualquer empresa que deseje fazê-lo está em má situação, incapaz de obter capital de bancos ou outras fontes. Em outro lugar, nós (Greenwald, Stiglitz e Weiss (1984)) fornecemos um modelo de seleção adversa simples (novamente, uma variante do modelo de seleção adversa padrão) mostrando que apenas as piores empresas irão, de fato, recorrer ao mercado de ações para levantar capital.

2. A tentativa de Keynes de explicar as flutuações econômicas apenas em termos de considerações de demanda não apenas colocou o dilema a que nos referimos anteriormente — por que as empresas não usam a política de preços para aumentar suas vendas — mas colocou outro problema: como poderia uma pequena economia aberta enfrentar problemas de desemprego keynesiano? Simplesmente alterando sua taxa de câmbio, ela poderia enfrentar uma demanda ilimitada por seus produtos.

Em nossa teoria, não há uma distinção clara entre demanda e oferta. As empresas estariam dispostas a produzir mais, se pudessem ter uma demanda garantida. Nesse sentido, a demanda está limitando a produção. As empresas não estão dispostas a produzir mais, dados os riscos associados à produção na ausência de uma demanda garantida. Nesse sentido, as empresas estão em sua curva de oferta.

Nossa teoria explica, assim, por que a quantidade de bens que as empresas estão dispostas a fornecer a qualquer salário real do produto esperado pode mudar ao longo do ciclo econômico.

Nossa teoria também pode explicar por que as empresas, ao definirem seus preços, podem tentar ter uma margem de lucro (mark-up) mais alta sobre os custos em períodos recessivos. Em mercados com competição imperfeita e informação imperfeita, as empresas devem recrutar clientes. Elas o fazem parcialmente usando políticas de preços. Elas enfrentam, assim, uma escolha (trade-off): preços mais baixos hoje levando a vendas futuras mais altas, lucros futuros mais altos, mas lucros correntes mais baixos. O preço que escolhem depende do custo implícito do capital (não da taxa de juros de mercado) e, na presença de racionamento de capital próprio, este pode ser mais alto em períodos recessivos.

3. Keynes argumentou que o determinante primário do nível de investimento, dado um conjunto de expectativas, era a taxa de juros. Embora sempre tenha havido alguma ambiguidade sobre se esta é a taxa de juros real ou nominal, o único sentido que os economistas modernos podem fazer disso é que deve ter sido a taxa de juros real. Mas as taxas de juros reais de mercado flutuaram relativamente pouco (até os anos 80). Uma boa teoria nunca deveria tomar uma constante (ou uma quase constante) como uma variável explicativa.

Em nossa teoria, a disponibilidade de crédito em certos momentos é o principal determinante do nível de investimento. É precisamente nesses momentos que a política monetária pode afetar o nível de atividade econômica. Em períodos recessivos, no entanto, os bancos podem estar dispostos a emprestar a qualquer "bom" prospecto à taxa de juros vigente, mas há uma escassez de tomadores de empréstimo dispostos. Em tais circunstâncias, a política monetária provavelmente será ineficaz.

A teoria keynesiana-neoclássica simplesmente não consegue explicar as flutuações de estoques, o fato de que os estoques servem para exacerbar em vez de amortecer as flutuações. Nossa teoria consegue. Novamente, o aumento no custo efetivo do capital — o resultado do racionamento de capital próprio e da diminuição na oferta de capital de giro — implica que as empresas desejarão diminuir seus estoques em períodos recessivos.

4. O mecanismo pelo qual as autoridades monetárias afetavam o nível de atividade econômica na análise keynesiana é implausível. Existem três etapas: (a) o governo toma ações que afetam a oferta de moeda; (b) dadas as funções de demanda dos indivíduos por moeda (uma função presumivelmente de taxas de juros e renda), as taxas de juros mudam; (c) como resultado das mudanças nas taxas de juros, o investimento muda.\footnote{Isto é obviamente uma supersimplificação. Em algumas variantes da teoria, a demanda por moeda depende apenas da renda e, portanto, dadas as rigidezes de preços, uma diminuição na oferta de moeda deve ser acompanhada por uma diminuição na renda. Nenhum mecanismo plausível pelo qual isso é efetuado foi apresentado. Em outras teorias, a demanda por investimento é uma função da renda futura esperada, que por sua vez é uma função da renda atual. As flutuações no investimento tornam-se então tanto uma consequência quanto uma causa das flutuações de renda. É difícil reconciliar tais modelos ingênuos de acelerador com o comportamento racional.}

Existem problemas com cada uma das etapas: embora o governo possa ser capaz de afetar a oferta de moeda externa (outside money), existem substitutos próximos da quase-moeda, pelo menos para fins de transações. Além disso, a moeda não é necessária para a maioria das transações, apenas o crédito. (É isso que torna esses modelos baseados na restrição de adiantamento de caixa (cash-in-advance) tão implausíveis.) E na medida em que a moeda é necessária para fins de transações, deve-se explicar por que isso acontece. Além disso, a relação entre transações e renda é tênue: muitas, talvez a maioria, das transações são trocas de ativos, e os tipos de mudanças econômicas associadas ao ciclo econômico são frequentemente acompanhados por mudanças na riqueza e, portanto, na distribuição de ativos.

Na medida em que a demanda por moeda é baseada em considerações de ativos, o que é relevante, é claro, não é a renda, mas a riqueza. E como existem títulos de curto prazo que são, exceto para fins de transações, substitutos perfeitos para a moeda, o custo de oportunidade relevante de manter moeda é a taxa de juros de curto prazo; mas se qualquer taxa de juros for relevante para o investimento, deve ser a taxa de juros real.\footnote{Não está claro se deve ser a taxa de juros real de longo ou curto prazo. Quando a questão é: quando um projeto deve ser empreendido, a taxa de juros real de curto prazo é presumivelmente relevante; quando a questão é: deve um projeto ser empreendido, é presumivelmente a taxa de juros real de longo prazo. Como a informação relevante para empreender um projeto (o conjunto de fornecedores, os preços pelos quais os fatores podem ser comprados, etc.) torna-se obsoleta tão rapidamente, em muitos casos, pelo menos, a questão colocada pelas empresas é mais a última do que a primeira.} Além disso, como o desenvolvimento recente de Contas de Gestão de Caixa (Cash Management Accounts) deixa claro, é perfeitamente viável fornecer "moeda" que rende juros, caso em que a única questão relevante enfrentada pelo indivíduo é a estrutura de vencimento da dívida que ele deseja manter.

Keynesianos mais recentes (por exemplo, Tobin) propuseram outro mecanismo pelo qual a política monetária afeta a atividade econômica: Na abordagem geral de portfólio, diferentes ativos (curto prazo, títulos de longo prazo) são vistos como substitutos imperfeitos, e mudanças na oferta relativa afetam diferentes taxas de juros e, em particular, o preço das ações. Isso pode ser criticado em vários fundamentos. Primeiro, as empresas não recorrem, em sua maioria, ao mercado de ações para levantar capital. Assim, o preço das ações não é diretamente relevante. Como podemos explicar as correlações observadas? Em nossa teoria, expectativas otimistas, digamos sobre vendas futuras, refletir-se-ão em um alto preço das ações (altos lucros futuros) e na disposição dos gestores em produzir. Há uma correlação, mas não causalidade.

Para colocar de outra forma, o que preocupa os gestores e os acionistas controladores não é o preço das ações hoje, mas o preço das ações quando eles forem vender suas ações. O preço atual pode ser uma boa previsão de preços futuros, mas os empresários têm mais probabilidade de basear seus julgamentos sobre projetos de investimento específicos não nos julgamentos de algum estranho relativamente mal informado, mas em suas próprias visões internas melhor informadas.

Em segundo lugar, em teoria, mudanças na estrutura de vencimento da dívida do governo não deveriam ter efeito no equilíbrio de mercado, desde que não existam consequências redistributivas significativas dessa mudança (e estas parecem implausíveis).\footnote{Isto não quer dizer que os indivíduos não levem em conta suas obrigações tributárias futuras; mas quais evidências existem que apoiam a hipótese de que eles levam essas obrigações tributárias futuras "totalmente" em conta? Não indicamos todos os pressupostos relevantes nos "Teoremas de Irrelevância", apenas o que consideramos ser o mais importante. Assim, por exemplo, os teoremas de irrelevância assumem que a tributação não é distorciva. Mas os impostos são distorcivos. Reduzir os impostos hoje e aumentar os impostos no futuro pode ter um efeito real no bem-estar; mas parece implausível que as diferenças nesses triângulos de distorção de Harberger, sendo cada um deles pequeno, possam explicar o fato de a política monetária ter qualquer efeito significativo.} Pois essas mudanças representam mudanças nas responsabilidades tributárias (estocásticas) futuras dos indivíduos. Os indivíduos, ao decidirem sobre suas carteiras ótimas, devem levar em conta outros aspectos dos riscos que enfrentam, incluindo riscos de salário e impostos; e se eles fizerem isso corretamente, não haverá efeitos nas taxas de juros reais. A abordagem de Tobin parece basear-se em um comportamento irracional. (Veja Stiglitz 1981)

Tobin pode objetar a isso em dois fundamentos: primeiro, que os indivíduos não incluem realmente suas responsabilidades tributárias estocásticas em sua análise de portfólio; e segundo, que nossa análise assumiu um mercado de capitais perfeito. Estamos inclinados a concordar; como resultado, também concordamos que o governo pode mudar as taxas de juros de mercado. Mas permanecemos não convencidos de que este seja o mecanismo (importante) pelo qual o governo controla a atividade econômica. Em vez disso, argumentaríamos que a política monetária do governo afeta a disposição dos bancos em emprestar (e os termos sob os quais eles estão dispostos a emprestar), e é através deste mecanismo que o investimento pode ser afetado.

\section*{Observações Finais}

\textbf{Sobre a metodologia.} As economias capitalistas são complicadas. Um modelo deve capturar suas características centrais, não reproduzi-las exatamente. As decisões de indivíduos e empresas hoje são baseadas em expectativas futuras e são afetadas por decisões passadas. Os indivíduos não têm previdência perfeita ou expectativas racionais em relação ao futuro. Os eventos que enfrentam frequentemente parecem ser únicos, e não há como formar um modelo estatístico prevendo a distribuição de probabilidade dos resultados. E há pouca evidência de que eles sequer tentem fazê-lo. Ao mesmo tempo, os indivíduos não são míopes. Eles não simplesmente assumem que o futuro é como o presente.

Os mercados não são perfeitos. Mas os mercados existem. Os preços se ajustam. Os salários caem na presença de desemprego massivo. Estes "fatos" impõem algumas decisões estratégicas importantes para o modelador: no futuro previsível, não é possível construir um modelo dinâmico que reflita adequadamente todos eles. Casos polares são mais fáceis de estudar. Deve-se assumir flexibilidade perfeita de salários ou preços ou nenhuma flexibilidade de salários ou preços? Expectativas racionais ou miopia? Qualquer conjunto de escolhas está aberto a críticas, mas igualmente, pode ser defendido como parte de uma estratégia de pesquisa de longo prazo.

Em nossa visão, as escolhas devem ser ditadas pelo fenômeno a ser estudado. O problema central em que estamos interessados é explicar o desemprego. Assim, começar a análise assumindo a limpeza do mercado é ignorar o que deve ser explicado.

Embora concordemos sobre a importância de entender os problemas de maximização dinâmica em que indivíduos e empresas estão engajados, ignorar as restrições importantes que enfrentam (por exemplo, no acesso aos mercados de capitais) resulta em modelos que têm pouca relevância. Suspeitamos que, em muitos casos, modelos míopes focando nas restrições são muito melhores do que modelos "racionais" que as ignoram. De fato, em alguns casos, pode-se mostrar que os modelos racionais com restrições parecem idênticos aos modelos míopes padrão (por exemplo, com comportamento impulsionado por regras, todos os lucros e nenhum dos salários poupados).\footnote{A escola de expectativas racionais é frequentemente creditada com a observação de que as políticas governamentais, se antecipadas, terão efeitos bastante diferentes dos pretendidos. Isso parece estar dando a eles mais crédito do que o devido. Nos anos 60, durante o período em que modelos dinâmicos estavam sendo tão ativamente investigados, houve uma quantidade considerável de trabalho analisando modelos de previdência perfeita; incluída nessas análises estava a análise das consequências de ações governamentais totalmente antecipadas. Notou-se, por exemplo, que um crédito tributário de investimento temporário diminuiria o investimento antes de sua data efetiva, causaria um surto de investimento subsequente, causaria outro surto de investimento antes de sua remoção e levaria a um investimento deprimido imediatamente após sua remoção. A contribuição das análises de expectativas racionais foi investigar essas questões em um ambiente estocástico. Embora esta tenha sido uma extensão importante dos modelos de previdência perfeita não estocásticos, as percepções básicas — pelo menos a observada acima — permanecem as mesmas.}

\textbf{Sobre a política.} Tem havido uma controvérsia de longa data sobre o que os governos devem fazer diante do desemprego: (a) nada; (b) encorajar reduções salariais; (c) usar política monetária; ou (d) aumentar os gastos governamentais. O sucesso da teoria keynesiana teve muito a ver com o fato de fornecer uma justificativa teórica para aqueles que desejavam seguir o quarto caminho. O sucesso da Nova Teoria Clássica teve muito a ver com o fato de ter fornecido uma justificativa teórica para aqueles que desejavam que o governo não fizesse nada.

Em nossa visão, a análise de Keynes estava basicamente correta. A política governamental pode afetar o resultado; em períodos recessivos, a política monetária provavelmente será de eficácia limitada; e os cortes salariais podem não ser eficazes.\footnote{Na teoria keynesiana, os cortes salariais reduzem a demanda agregada. Na teoria keynesiana mais moderna, onde o consumo é baseado na renda permanente, tais cortes salariais podem ter um efeito insignificante na demanda. Nossa teoria fornece uma explicação de por que os cortes salariais podem ter um efeito significativo: os mercados de capitais imperfeitos resultam em alguns indivíduos tendo que reduzir seu consumo. Por outro lado, existem circunstâncias em nossa teoria onde um corte salarial seria eficaz: quando cada empresa opta por não reduzir seu salário, dados os salários pagos por outras empresas, uma mudança salarial coordenada pode aumentar a demanda por trabalho. Nossa teoria sugere, no entanto, que existem outras circunstâncias onde a redução dos salários reais (abaixo do salário de eficiência) resultaria na verdade em uma redução na demanda por trabalho. Na medida em que salários mais baixos levam a preços mais baixos, reduções salariais podem ter efeitos deletérios futuros, ao reduzir o capital de giro disponível para as empresas e ao torná-las mais relutantes em produzir, se extrapolarem os declínios atuais nos preços para continuarem no futuro.}

\textbf{Sobre a eficiência da economia de mercado.} Concordamos com Keynes que o desemprego é um problema real que as economias capitalistas enfrentam. Embora meio século de experiência possa nos tornar menos otimistas sobre a capacidade do governo de eliminar as flutuações econômicas, meio século de experiência com formas alternativas de organização econômica nos tornou ainda menos otimistas sobre a capacidade dessas alternativas de fornecer a base de um sistema mais eficiente de alocação de recursos. Como nas roupas novas do imperador, podemos não ser capazes de ver a mão invisível porque ela não está lá; ou talvez com mais precisão, porque ela está tão invisível que não vemos quão paralisada ela está. O desemprego é apenas a pior manifestação de falhas de mercado generalizadas que surgem na presença de informação imperfeita e mercados incompletos. Mas se a mão invisível do mercado está paralisada, a mão visível do governo pode ser muito pior. Voltaire estava errado: não vivemos no melhor dos mundos possíveis. Vivemos em um mundo imperfeito. E devemos aprender a viver com essas imperfeições. A intervenção governamental limitada — corrigindo as piores manifestações das falhas de mercado, incluindo o desemprego massivo — pode, afinal, melhorar a eficiência da economia de mercado. No final, Keynes e as políticas keynesianas são vindicados.

\newpage

\section*{Referências}

Akerlof, G. [1970], "The Market for Lemons: Qualitative Uncertainty and the Market Mechanism," \textit{Quarterly Journal of Economics}, Vol. 84, pp. 288-300.

Akerlof, G. and J. Yellen, "A Near Rational Model of the Business Cycle with Wage and Price Inertia," \textit{Quarterly Journal of Economics}, 1985, pp. 832-38.

Barro, R. J., and Grossman, H. I. (1971). "A general disequilibrium model of income and employment." \textit{American Economic Review} 61:82-93.

Blinder, A. and J. E. Stiglitz. "Money, Credit Constraints and Economic Activity," \textit{American Economic Review}, May 1983, pp. 297-302.

Greenwald, B. [1979]. \textit{Adverse Selection in the Labor Market}, New York, Garland Press.

Greenwald, B. and J. E. Stiglitz. "Externalities in Economies with Imperfect Information and Incomplete Markets," \textit{Quarterly Journal of Economics}, May 1986.

Greenwald, B., Stiglitz, J.E. [1986], "Information, Finance Constraints and Business Fluctuations," \textit{Proceedings of the Taiwan Conference on Monetary Theory}, Taipei, Chung-Hua Institute.

Greenwald, B. and J. E. Stiglitz. "Money, Imperfect Information, and Economic Fluctuations," artigo apresentado na \textit{Conference on Monetary Theory} em Taipei, Janeiro, 1986, e publicado em \textit{Proceedings} (Oxford University Press).

Greenwald, B., Stiglitz, J.E. and Weiss, A. M. [1984], "Informational Imperfections and Macroeconomic Fluctuations," \textit{American Economic Review, Papers and Proceedings}, Vol. 74, pp. 194-99.

Hansen, B. (1951). \textit{A study in the theory of inflation}. Allen and Unwin, London.

Kalecki, M. [1939], \textit{Essays in the Theory of Economic Fluctuations}, New York, Russell and Russell.

Myers, S. C. and Majluf, N. S. [1984], "Corporate Financing and Investment Decisions When Firms Have Information that Investors Do Not," \textit{Journal of Financial Economics}, Vol. 11, pp. 187-221.

Neary, P., and J. E. Stiglitz. "Towards a Reconstruction of Keynesian Economics: Expectations and Constrained Equilibria," \textit{Quarterly Journal of Economics}, Supplement 1983, pp. 199-228 (apresentado no encontro de Atenas da Econometric Society, Setembro de 1979).

Phelps, E. S. and Winter, S. G. [1970], "Optimal Price Policy Under Atomistic Competition" em E. S. Phelps, ed., \textit{Microeconomic Foundations of Employment and Inflation Theory}, New York. W. W. Norton and Co.

Solow, R. , and J. E. Stiglitz. "Output, Employment and Wages in the Short Run," \textit{Quarterly Journal of Economics}, Vol. 82, November 1968, pp. 537-560.

Stiglitz, J. E. (1986b) "The Causes and Consequences of the Dependence of Quality on Price," Hoover, mimeo.

J. E. Stiglitz (1986a), "Theories of Wage Rigidities," artigo apresentado na conferência \textit{Keynes' Economic Legacy}, University of Delaware, Dezembro de 1983, em \textit{Keynes' Economic Legacy}, Praeger Publishers, 1986.

Stiglitz, J. E. [1985], "Credit Markets and the Control of Capital," \textit{Journal of Money, Credit and Banking}, Vol. 17, pp. 133-152.

J. E. Stiglitz, "Price Rigidities and Market Structure," \textit{American Economic Review}, Vol. 74, No. 2, May 1984, pp. 350-56 (artigo apresentado na American Economic Association, Dezembro de 1983).

J. E. Stiglitz (1983), "Competition and the Number of Firms in a Market: Are Duopolies More Competitive than Atomistic Markets?", \textit{Journal of Political Economy}.

J. E. Stiglitz (1983), "On the Relevance of Irrelevance of Public Financial Policy: Indexation, Price Rigidities and Optimal Monetary Policy," artigo apresentado na \textit{Conference at Rio de Janeiro}, Dezembro de 1981, publicado em \textit{proceedings, Inflation, Debt and Indexation}, R. Dornbusch and M. Simonsen, eds., M.I.T. Press, 1983, pp. 183-222.

Stiglitz, J. E. (1981) "The Relevance or Irrelevance of Government Financial Policies," NBER working paper (publicado em M. Boskin, ed. \textit{Proceedings of IEA Roundtable on Budget Deficits}).

J. E. Stiglitz and A. Weiss, "Credit Rationing and Collateral," artigo apresentado na conferência patrocinada pelo CEPR, Oxford, Setembro de 1985.

Stiglitz, Joseph and A. Weiss. "Incentive Effects of Terminations: Applications to the Credit and Labor Markets," \textit{American Economic Review}, 73, (1983): 912-927.

Stiglitz, J. E. and Weiss, A. [1981], "Credit Rationing in Markets with Imperfect Information," \textit{American Economic Review}, Vol. 71, pp. 393-410.

\end{document}